\chapter{Introdução}
\label{cap:introducao}

Apresente o tema, o problema de pesquisa, a justificativa, os objetivos e a
organização do trabalho. Delimite o escopo para que as perguntas que orientam
o estudo sejam verificáveis.

\section{Motivação e aplicações}
\begin{frame}{Por que o oscilador harmônico?}
  \begin{columns}[T,onlytextwidth]
    \begin{column}{0.54\textwidth}
      \begin{itemize}
        \item Perto de um mínimo local com \alert{$V''(x_0)>0$},
              o potencial é aproximadamente parabólico:
      \end{itemize}
      \begin{equacao}
        \[
          V(x) \approx V(x_0)
          + \tfrac12 V''(x_0)\,(x-x_0)^2
        \]
      \end{equacao}
      \begin{itemize}
        \item Por isso ele descreve \alert{moléculas}, \alert{sólidos},
              \alert{campos} e \alert{circuitos}.
        \item É um dos poucos modelos quânticos com solução
              \alert{analítica exata}.
      \end{itemize}
    \end{column}
    \begin{column}{0.42\textwidth}
      \begin{destaque}{Onde ele aparece}
        \begin{itemize}
          \item Vibrações moleculares (espectroscopia IV)
          \item Fônons em cristais
          \item Íons aprisionados e qubits
          \item Ressonadores nanomecânicos
        \end{itemize}
      \end{destaque}
    \end{column}
  \end{columns}
\end{frame}

\begin{frame}{Onde isso aparece}
  \begin{columns}[T,onlytextwidth]
    \begin{column}{0.49\textwidth}
      \begin{destaque}{Moléculas diatômicas}
        Vibrações em torno do equilíbrio. As transições
        $\Delta n = \pm 1$ explicam as bandas no infravermelho.
      \end{destaque}
      \begin{destaque}[fisicaamber]{Íons aprisionados}
        O movimento de um íon em um potencial efetivo forma um
        oscilador; seus estados são usados em computação quântica.
      \end{destaque}
    \end{column}
    \begin{column}{0.49\textwidth}
      \begin{destaque}[fisicaviolet]{Circuitos supercondutores}
        Junções Josephson formam um oscilador anarmônico cujos dois
        primeiros níveis definem um qubit.
      \end{destaque}
      \begin{destaque}[fisicagreen]{Nanomecânica}
        Osciladores mecânicos levados ao estado fundamental testam
        os limites da mecânica quântica.
      \end{destaque}
    \end{column}
  \end{columns}
\end{frame}


\section{Objetivos}
\label{sec:objetivos}

O objetivo geral deste exemplo é analisar o oscilador harmônico quântico.
Como objetivos específicos, propõem-se:
\begin{itemize}
  \item formular o problema pela equação de Schrödinger;
  \item determinar e interpretar os níveis de energia;
  \item comparar as previsões quântica e clássica.
\end{itemize}

\section{Organização do trabalho}

O Capítulo~\ref{cap:fundamentacao} apresenta o referencial teórico. O
Capítulo~\ref{cap:metodologia} descreve o método, e o
Capítulo~\ref{cap:resultados} reúne os resultados e sua discussão. O
Capítulo~\ref{cap:conclusao} encerra o texto.
